% Page 043 — page imprimée 39
% Contenu seul : le préambule, l'en-tête et la pagination sont gérés par meyrueis.tex

\section*{MEYRUEIS}

\begin{center}
\textbf{Bonnes pages des souvenirs de Jean Gounelle (1902-1994)}

\vspace{0.4em}
\textbf{Pasteur de Meyrueis de 1932 à 1935}
\end{center}

Avec Meyrueis le jeune ménage pastoral que nous formions, Irène et moi, abordait un
pays plus civilisé que le Pompidou. Le territoire paroissial qui, désormais, allait nous
incomber comportait une petite ville et trois vallées.

Meyrueis constituait alors, comme aujourd’hui encore, un centre avec docteur,
dentiste, PTT, pharmacien, garage, Caisse d’Epargne, Banque, et beaucoup de
magasins etc. pourvu en outre de nombreux hôtels, dont certains luxueux. On sentait
la petite cité touristique, soignée, attirante, se suffisant à elle même, susceptible
d’accueillir une clientèle exigeante.

La population de ce bourg se rattachait, en grande partie, à l’Eglise Catholique. Les
protestants, majoritaires au 19\textsuperscript{e} siècle, ayant les premiers et assez massivement quitté
la ville, les Caussenards venaient combler les vides ainsi créés. L’école laïque ne
conservait qu’une influence secondaire par rapport à l’école libre, tenue par de jeunes
frères, instituteurs en soutane, très dynamiques et animés d’un grand zèle. De plus s’y
trouvait implantée une maison de « sœurs des pauvres » Ces femmes, d’un
dévouement sans borne, travaillaient dans ce milieu, appréciées de tous.
Notre ministère où la partie médicale avait tenu un si grand rôle au Pompidou, ne
pouvait qu’en être profondément modifié. Ce côté cependant, n’en fut pas pour
autant négligeable. !!

Notre action s’engagea surtout dans les trois vallées, entièrement protestantes elles,
celle de la Jonte, de la Brèze et du Bétuzon, trois ruisseaux dont les noms poétiques
chantaient comme le bruissement de leurs eaux. Ces trois vallées se rejoignaient à
l’endroit même où n’avait pu que surgir, dans la nuit des temps, un groupement
humain, protégé, cohérent et naturel.
Je ne vois pas d’autre explication à la naissance et à l’origine de Meyrueis.

Pour rester dans les considérations générales et m’en tenir à ce qui constituait ma
charge en ce lieu, disons que la paroisse protestante, dans les années 1930, avait un
double visage. Celui de l’hiver dans l’axe de ces trois vallées où se trouvaient nos
paysans, beaucoup plus à l’aise que ceux du Pompidou, à cause de leur lait de brebis
exploité par les caves de Roquefort. Celui de l’été. De Juin à Octobre, l’auditoire au
culte, nombreux et divers, se composait surtout de fidèles issus de grandes paroisses
citadines. Le Château d’Ayres, en particulier, dont le propriétaire faisait partie du
Conseil Presbytéral, constituait un pôle d’attraction exceptionnel.

Nos familles bourgeoises protestantes de Paris et d’ailleurs venaient volontiers y
passer l’été et revenaient ensuite. Elles animaient, nos « ventes », nos fêtes d’Eglise
organisaient au Temple de nombreuses rencontres et parfois de magnifiques concerts.
Le pasteur devait donc s’adapter à ce double aspect de l’Eglise locale, et ne pas
négliger les uns aux dépens des autres. Il convient également de se souvenir qu’en ces
